% Public content API: metadata, structure, and reusable study-note blocks.
\makeatletter
\newcommand{\StudySubject}{}
\newcommand{\StudyCode}{}
\newcommand{\StudyYear}{}
\newcommand{\StudyAuthor}{}
\newcommand{\subject}[1]{\gdef\StudySubject{#1}}
\newcommand{\subjectcode}[1]{\gdef\StudyCode{#1}}
\newcommand{\academicyear}[1]{\gdef\StudyYear{#1}}
\newcommand{\authorname}[1]{\gdef\StudyAuthor{#1}}

% Minimal title page; optional metadata rows disappear when left empty.
\newcommand{\makestudytitlepage}{%
  \begin{titlepage}
    \thispagestyle{empty}
    \vspace*{0.19\textheight}
    {\color{PrimaryColor}\rule{\textwidth}{1.2pt}}\par
    \vspace{1.3cm}
    {\Huge\bfseries\color{PrimaryColor}\StudySubject\par}
    \vspace{0.35cm}
    {\Large Study Notes\par}
    \ifx\StudyCode\@empty\else
      \vspace{0.35cm}{\large\StudyCode\par}
    \fi
    \ifx\StudyYear\@empty\else
      \vspace{0.2cm}{\large Academic Year \StudyYear\par}
    \fi
    \vfill
    \ifx\StudyAuthor\@empty\else
      {\large\StudyAuthor\par}
    \fi
    \vspace*{1.1cm}
    {\color{PrimaryColor}\rule{\textwidth}{0.5pt}}\par
  \end{titlepage}%
}

% Numbered chapter heading; using refstepcounter keeps section and question
% counters chapter-aware without requiring manually entered chapter numbers.
\newcommand{\chaptertopic}[1]{%
  \clearpage
  \refstepcounter{chapter}%
  \setcounter{section}{0}%
  \setcounter{subsection}{0}%
  \markboth{\StudySubject}{Chapter \thechapter\ -- #1}%
  \addcontentsline{toc}{chapter}{Chapter \thechapter\enspace #1}%
  \thispagestyle{plain}%
  {\color{PrimaryColor}\large\bfseries Chapter \thechapter\par}
  \vspace{0.4em}
  {\Huge\bfseries #1\par}
  \vspace{0.7em}
  {\color{RuleColor}\rule{\textwidth}{0.6pt}\par}
  \vspace{1.2em}
}

% Chapter-scoped counters make labels and printed numbers agree (e.g. 2.3).
\newcounter{studyquestion}[chapter]
\renewcommand{\thestudyquestion}{\thechapter.\arabic{studyquestion}}
\newcounter{studydefinition}[chapter]
\renewcommand{\thestudydefinition}{\thechapter.\arabic{studydefinition}}

% A definition box uses a shared style and accepts an optional label.
\newtcolorbox{StudyDefinitionBox}[1]{
  definitionbox,
  title={#1}
}
\NewDocumentCommand{\definition}{o m +m}{%
  \refstepcounter{studydefinition}%
  \IfNoValueF{#1}{\label{#1}}%
  \begin{StudyDefinitionBox}{#2}#3\end{StudyDefinitionBox}%
}

% Questions are list items and must be placed inside questionlist.
% Optional labels use ordinary LaTeX references, e.g. \q[q:bpm]{...}.
\newenvironment{questionlist}{%
  \begin{list}{}{%
    \setlength{\leftmargin}{3.7em}%
    \setlength{\labelwidth}{3.3em}%
    \setlength{\labelsep}{0.4em}%
    \setlength{\itemsep}{0.7\baselineskip}%
    \setlength{\parsep}{0.2\baselineskip}%
    \setlength{\topsep}{0.5\baselineskip}%
    \setlength{\listparindent}{0pt}%
  }%
}{\end{list}}
\NewDocumentCommand{\q}{o +m}{%
  \refstepcounter{studyquestion}%
  \item[\textbf{\color{PrimaryColor}Q\thestudyquestion.}]%
  \IfNoValueF{#1}{\label{#1}}%
  #2\par
}
\NewDocumentCommand{\answer}{+m}{%
  \par\smallskip
  {\small\color{TextMuted}\textbf{Answer:} #1\par}%
  \smallskip
}
\NewDocumentCommand{\qa}{o +m +m}{%
  \IfNoValueTF{#1}{\q{#2}}{\q[#1]{#2}}%
  \answer{#3}%
}
\NewDocumentCommand{\examquestion}{o +m}{%
  \refstepcounter{studyquestion}%
  \item[\textbf{\color{PrimaryColor}Exam Q\thestudyquestion.}]%
  \IfNoValueF{#1}{\label{#1}}%
  \begin{tcolorbox}[
    studybox,
    colback=ExamBackground,
    colframe=PrimaryColor!65,
    borderline west={1.3pt}{0pt}{PrimaryColor},
    before skip=0.35\baselineskip,
    after skip=0.3\baselineskip
  ]#2\end{tcolorbox}%
}
\NewDocumentCommand{\examqa}{o +m +m}{%
  \IfNoValueTF{#1}{\examquestion{#2}}{\examquestion[#1]{#2}}%
  \answer{#3}%
}

% Other reusable blocks share tcolorbox geometry, but have distinct roles.
\NewDocumentCommand{\keyconcept}{m +m}{%
  \begin{tcolorbox}[conceptbox,title={Key Concept\quad #1}]#2\end{tcolorbox}%
}
\NewDocumentCommand{\important}{+m}{%
  \begin{tcolorbox}[importantbox,title={Important}]#1\end{tcolorbox}%
}
\NewDocumentCommand{\note}{+m}{%
  \begin{tcolorbox}[notebox,title={Note}]#1\end{tcolorbox}%
}
\NewDocumentCommand{\example}{m +m}{%
  \begin{tcolorbox}[examplebox,title={Example\quad #1}]#2\end{tcolorbox}%
}
\newcommand{\term}[1]{\textcolor{PrimaryColor}{\textbf{#1}}}
\newtcolorbox{summarybox}{summarybox}
\newenvironment{learningobjectives}{%
  \begin{tcolorbox}[studybox,colback=LightBackground,colframe=RuleColor,title={Learning Objectives}]%
  \begin{enumerate}[label=\textcolor{PrimaryColor}{\arabic*.},leftmargin=*]
}{%
  \end{enumerate}%
  \end{tcolorbox}%
}

% Select the listings lexer by name. C# is entered as C\# in LaTeX source.
\newcommand{\StudyCodeLanguage}[1]{%
  \IfStrEqCase{#1}{%
    {C\#}{\lstset{language=CSharp}}%
    {CSharp}{\lstset{language=CSharp}}%
    {JavaScript}{\lstset{language=JavaScript}}%
    {TypeScript}{\lstset{language=TypeScript}}%
    {Python}{\lstset{language=Python}}%
    {SQL}{\lstset{language=SQL}}%
    {HTML}{\lstset{language=HTML}}%
    {CSS}{\lstset{language=CSS}}%
    {Bash}{\lstset{language=bash}}%
  }[\lstset{language=#1}]%
}
\lstnewenvironment{codeblock}[1]{%
  \lstset{style=StudyCode}%
  \StudyCodeLanguage{#1}%
}{}
\makeatother
